% Generated draft: requires mathematical and source review.
\documentclass[11pt]{article}
\usepackage[margin=25mm]{geometry}
\usepackage{fontspec}
\setmainfont{Latin Modern Roman}
\usepackage{amsmath,amssymb,mathtools,unicode-math}
\setmathfont{Latin Modern Math}
\usepackage{enumitem,needspace,xurl,hyperref}
\setlength{\parindent}{0pt}
\setlength{\parskip}{5pt}
\setlength{\emergencystretch}{3em}
\newcommand{\N}{\mathbb{N}}
\newcommand{\Z}{\mathbb{Z}}
\newcommand{\Q}{\mathbb{Q}}
\newcommand{\R}{\mathbb{R}}
\newcommand{\C}{\mathbb{C}}
\newcommand{\F}{\mathbb{F}}
\newcommand{\A}{\mathbb{A}}
\newcommand{\PP}{\mathbb{P}}
\newcommand{\E}{\mathbb{E}}
\newcommand{\eps}{\varepsilon}
\newcommand{\dd}{\,\mathrm{d}}
\begin{document}
\section*{UnsolvedMath ID: 11000149}
\textbf{Source problem code:} catalogue entry
\subsection{Definitions and mathematical statement}
Let \(F\) be the source's fiber surface with one boundary component, let \(t_\partial\) be its positive boundary Dehn twist, and consider positive factorizations
\[t_\partial=t_{\alpha_1}t_{\alpha_2}\cdots t_{\alpha_k}\quad\text{in }\operatorname{Mod}(F),\]
where each \(t_{\alpha_i}\) is a right-handed twist. The main question is whether every such relation is one of the source's seven relations (1)--(7), up to its positive equivalence in the relevant Artin group. Those seven relations, surface/genus data, Artin group and equivalence moves are absent, so the proposed classification cannot be expanded into explicit cases. In particular positive Artin equivalence is not silently identified with arbitrary mapping-class equality or Hurwitz equivalence.

The symplectic motivation is a Lefschetz fibration \(p:M^4\to D\) over a disk, whose monodromy is the ordered product of positive twists in vanishing cycles \(\alpha_i\), associated with a geometric basis of \(\pi_1(D\setminus\{x_1,\ldots,x_k\},p_0)\). The \(x_i\) are critical values and \(p_0\in\partial D\). The fiber is assumed to have one boundary component and boundary monodromy \(t_\partial\). Any proposed classification should specify the allowed choices and their equivalences, rather than infer a fibration classification from the word identity alone.

The source reports [12] that for a generic Lefschetz pencil of plane curves of any fixed degree, one can choose a geometric basis so the vanishing cycles form a bouquet; it states that pairs then have intersection zero when tangent or one otherwise. It asks whether something similar holds for generic curve pencils on arbitrary algebraic surfaces, or more generally symplectic Lefschetz fibrations without reducible fibers, and proposes the restricted version of the prior Smith question in which
\[i(\alpha_i,\alpha_j)\in\{0,1\}\quad(i\ne j).\]
Here \(i\) denotes geometric intersection number; the source's bouquet/tangency convention requires clarification. Such curves give the commuting or length-three braid relations for their twists, explaining the indicated geometric homomorphism from an Artin group. The exact additional meaning of ``small Artin group'' is not supplied. The restricted question and the general seven-relation completeness question are not asserted to be equivalent.

There is also an extraction/formulation issue in the plane-pencil report: it prints a geometric basis of \(\pi_1(\mathbb CP^2\setminus\{x_1,\ldots,x_k\})\), while the \(x_i\) are described as critical values of a pencil, whose usual parameter base is a curve. The displayed projective-space index, basepoint and treatment of pencil basepoints/fiber boundary components must be checked. No silent \(\mathbb CP^2\)-to-\(\mathbb CP^1\) replacement or one-boundary normalization is made.
\subsection{Short English statement}
Do seven specified positive Artin-type relations exhaust boundary-twist factorizations? The source also proposes a restricted Lefschetz-fibration question with pairwise vanishing-cycle intersection zero or one, but the relation list and equivalence rules are missing.
\subsection{Sources}
\begin{itemize}
\item Local UnsolvedMath registry, permanent ID \(11000149\), exact source code \texttt{catalogue entry}.
\item Farb, \emph{Problems on Mapping Class Groups and Related Topics}, unnumbered Question, PDF page 133; references to relations (1)--(7) and the earlier Smith question: \url{https://www.math.uchicago.edu/~farb/papers/mcgbook.pdf}. Registry access 2026-07-29; pdf-font-structure extraction.
\item Internal reference [12] is credited for the plane-pencil bouquet result; its full bibliographic entry and author attribution are absent.
\end{itemize}
\end{document}
